\scanpage{86}


Si $b_1$ \deleted{\unsure{en}} engendre $L_1$ (supposé inv.) alors les $b_1\otimes u_2$ engendrent $L_1\otimes L_2$, donc \deleted{\unsure{il}} $K=\operatorname{Im}(\sigma_1\otimes\sigma_2-{\rm id}_{E_1\otimes E_2})$ contient $L_1\otimes L_2$, plus précisément il est égal à

\[ L_1\otimes L_2+O\cdot\xi \]

où

\[ \xi=\overline{x_1\otimes x_2}-\unsure{x_2|x_1}\otimes x_2\equiv\deleted{b_2}b_1\oplus b_2\quad(\text{mod }L_1\otimes L_2) \]

\inserted{non nul mod. $L_1\otimes L_2$}

Si $L_1$ est de rang $1$, $L_2$ de rang $n$, $\operatorname{Im}(\sigma_1\otimes\sigma_2-{\rm id})$ est donc de rang \deleted{$1+n$,} égal à

\[ \operatorname{rang} L_1\otimes L_2+1=n+1 \]

tandis que

\[ \begin{aligned}\operatorname{rang}(\operatorname{Ker} *)&=\underset{\substack{\Vert\\\operatorname{rg} E_1\,\operatorname{rg} E_2\\\Vert\\2\cdot(n+1)}}{\operatorname{rang}(E_1\otimes E_2)}-\underset{\substack{\Vert\\1+\operatorname{rang}(L_1\otimes L_2)\\\Vert\\1+n}}{\operatorname{rang} E_1*E_2}\\ &=n+1\end{aligned} \]

d'où l'égalité.

On trouve donc

\uline{Proposition} : Supposons \deleted{\unsure{$E$}} $L_1,L_2$ de \inserted{loc. libres de} \deleted{type} fini.

Pour que $E_1\otimes E_2\mathbin{\unsure{-|\to}}E_1*E_2$ induise

\[ {E_1\otimes E_2}_{/\sigma_1\otimes\sigma_2}\xrightarrow{\sim}E_1*E_2, \]

il f. et s. que \deleted{\unsure{à son}} \deleted{tout point} \deleted{\unsure{s}} de $X_0=V(2.1_{\unsure{x|X}})$, $T_{10}$ \uline{ou} $T_{20}$ engendre $L_{10}$ resp. $L_{20}$. (Si $L_1,L_2$ sont tous les deux inversibles, et si $E_1,E_2$ munis de scindages, d'où $b_1\in\Gamma L_1,b_2\in\Gamma L_2$, cette condition équivaut à

\[ 2\,\mathcal O_X+b_1L_1^\vee+b_2 L_2^{\unsure{\vee|*}}=\underline O_X\quad\Bigr\} \]
% archive sheet number: 85
